\labsection{8}{Turn data-management results into visual evidence}{sec:lab08}

\begin{labproject}{Lab 8b visualization workflow}
\textbf{Coverage.} Chapter~\ref{chap:visualization} and the cleaned-data workflow from Lab 7.\\
\textbf{Goal.} Improve the Lab 7 report with at least two different plot types, then use one built-in R dataset to tell a data story.

\textbf{Part A --- Add visual evidence to the external-file report.}

The source requires at least two different plot/chart/graph types using the same external files as Lab 7. Because those external datasets were not supplied with this notes build, use the actual cleaned objects from your Lab 7 environment and choose plots that match the columns you have verified.

Example patterns from the plot families taught in Lab 8a:

\begin{lstlisting}[style=dsr]
# Example: categorical count bar chart
sex_counts <- table(titanic_cleaned$sex)
barplot(
  sex_counts,
  main = "Passenger count by sex",
  xlab = "Sex",
  ylab = "Count"
)

# Example: numeric distribution
hist(
  titanic_cleaned$fare,
  main = "Fare distribution",
  xlab = "Fare"
)
\end{lstlisting}

Do not state an insight until you have inspected the actual resulting counts/plot in your class dataset.

\textbf{Part B --- Storytelling with a built-in dataset.}

The handout allows any built-in dataset and gives \texttt{BJsales} as an example. Since \texttt{mtcars} is already used throughout Lab 8a, it is a natural option for a coherent practice story:

\begin{lstlisting}[style=dsr]
data(mtcars)

boxplot(
  mpg ~ cyl,
  data = mtcars,
  xlab = "Number of cylinders",
  ylab = "Miles per gallon",
  main = "Mileage by cylinder group"
)

plot(
  mtcars$wt,
  mtcars$mpg,
  xlab = "Weight",
  ylab = "Miles per gallon",
  main = "Weight versus mileage"
)
\end{lstlisting}

Use the displayed patterns to write the story in this order:
\begin{enumerate}
  \item Identify what each axis or group represents.
  \item Describe the strongest visible pattern without overstating causation.
  \item Connect the plot to a practical question a manager or analyst might ask.
  \item If a second plot supports the same observation, explain how the evidence complements the first plot.
\end{enumerate}
\end{labproject}

\begin{labdeliverable}
Submit the improved Lab 7 report with at least two different visualization types and a separate built-in-dataset story supported by plots produced from the Lab 8a functions.
\end{labdeliverable}
